\documentclass[11pt]{article}
\usepackage{icmlko}
\begin{document}

\title{두 해 만에 \texttt{t} 를 처음 썼다 --- 최근순은 구멍이고 연중 위치는
실력이다}
\author{월드모델 실험실}
\date{노트 447 · 2026년 8월 2일}
\maketitle

\begin{abstract}
노트 446 이 두 가지를 동시에 남겼다. \textbf{챔피언이 \texttt{t} 를 아예
안 쓴다}는 것과, \textbf{그렇다고 날짜를 그냥 넣으면 안 된다}는 것 ---
최근순 축은 날 판을 $+0.0186$ 올리면서 날짜 통제 판을 $-0.0294$ 내린다.
그래서 이번엔 \textbf{노트 446 이 만든 가드로 후보를 걸렀다.}

\begin{center}
\begin{tabular}{lrr}
\hline
후보 & 판 (날것) & 날짜 통제 판 \\
\hline
최근순 (노트 446) & $+0.0186$ & $-0.0294$ \\
\textbf{계절} (연중 위치) & $\mathbf{+0.0045}$ & $\mathbf{+0.0041}$ \\
밀도 (그 시기 출시 수) & $-0.0280$ & $-0.0407$ \\
간격 (직전 출시와의 거리) & $-0.0087$ & $-0.0122$ \\
셋 다 & $-0.0209$ & $-0.0358$ \\
\hline
\end{tabular}
\end{center}

\textbf{넷 중 하나만 통과한다.} 연중 위치는 유보 안에서 날짜에 \emph{단조가
아니라서}(2025$\sim$2026 을 두 번 돈다) 라벨의 \emph{쌓임}을 못 읽는다.
그래서 구멍을 안 판다.

관문에서도 둘 다 섰다.

\begin{center}
\begin{tabular}{lrrr}
\hline
짝 12뽑기 & 차 & SD & 부호 \\
\hline
판 & $+0.0041$ & $0.0042$ & $9/12$ \\
\textbf{날짜 통제 판} & $\mathbf{+0.0046}$ & $0.0038$ & $\mathbf{11/12}$ \\
\hline
\end{tabular}
\end{center}

\textbf{날짜 통제 쪽이 더 크고 더 고르다} --- 최근순과 정반대다.
\textbf{채택.}
\end{abstract}

\begin{figure}[t]
\centering
\includegraphics[width=\textwidth]{figs/f1.pdf}
\caption{\textbf{가} --- 후보 다섯을 두 판 위에 놓았다. 초록 사분면(둘 다
오름)에 든 것은 계절 하나다. 최근순은 오른쪽 아래 --- 날것만 오르는
자리다. \textbf{나} --- 계절의 짝 12뽑기. 오차 막대는 짝 SD.}
\end{figure}

\section{애니 하나가 끌고 간다}

도메인별로 보면 값이 고르게 퍼지지 않는다. \textbf{애니가 $+0.0264$ 로
$12/12$} 이고 나머지는 대부분 잡음 크기다(만화 $-0.0117$ 이 제일 큰
손해). 그런데 애니는 전이에서도 같은 일을 한다 --- LODO 에서
\textbf{$-0.0713$ 에서 $+0.1017$ 로 부호를 바꾼다}.

애니가 계절을 크게 타는 것은 그 자체로 말이 된다. \emph{분기 애니}는
1 · 4 · 7 · 10월에 몰려 나오고, 같은 분기 안에서 서로를 갉아먹는다.
연중 위치가 그 구조를 그대로 적는다.

\section{전이 쪽 확인}

노트 420 의 거부권(\emph{전이를 확실히 잃는 변경은 판이 통과해도 안
넣는다})을 걸어 봤다. LODO 에서 문턱($0.0108$)을 넘은 것이 넷뿐이라
노트 439 로 \textbf{잴 수 없음}이고, 그 넷은 \textbf{오름 셋 · 내림
하나}에 평균 $+0.0183$ 이다. 거부권은 \textbf{안 걸린다}.

\section{구현과 가드}

\texttt{F18\_bagboost.SEASON} 으로 넣었다. \texttt{fit} 과
\texttt{predict} 이 각각 \texttt{\_season(t)} 로 (값, 표시자) 두 열을
붙인다 --- \textbf{두 해 동안 모든 \texttt{predict} 가 받기만 하고 안
쓰던 인자를 처음 쓴다}.

가드를 두 개 걸었다.

\begin{itemize}
\item \textbf{정렬}(노트 414 의 교훈) --- 관문은 \texttt{d.yr} 로 쟀고
구현은 \texttt{years(t)} 로 편다. 열한 도메인 전부 최대 차
$0.00\mathrm{e}{+}00$.
\item \textbf{재현} --- 구현을 켜고 끈 판이 선별 때 수치와 소수점 넷째
자리까지 같다($+0.4764 \to +0.4809$, 통제 $+0.5158 \to +0.5199$).
\end{itemize}

되돌림 조항(노트 432): 되돌릴 상태의 이름은 \textbf{``계절 축 없는
36축 F18''}.

\section*{남은 것}

\begin{itemize}
\item \textbf{계절을 도메인마다 다르게.} 지금은 연중 위치 한 열을 열한
도메인이 나눠 쓴다. 애니만 크게 타는데, 애니의 분기 구조를 직접 적으면
(1 · 4 · 7 · 10월까지의 거리) 더 버나.
\item \textbf{달력 축 넷과의 겹말.} 이미 있는 달력 축과 얼마나 겹치는지
안 봤다. 겹말이면 둘 중 하나를 뗄 수 있다.
\item \textbf{밀도가 왜 그렇게 나쁜가}($-0.0407$). 경쟁 밀도는 말이 되는
재료인데 지금 재는 법(같은 도메인 $\pm 0.08$년 출시 수)이 유보에서
학습과 다른 뜻이 되는 듯하다.
\item LODO 씨앗 셋$\to$아홉 --- 여덟 노트째 미룸.
\end{itemize}

\end{document}
